% Source: https://artofproblemsolving.com/community/c3h1569382_even_numbers
% Stable identity: aops-topic-1569382-q01
Pair each positive integer $n$ with the even number $2n$:
\[
1\leftrightarrow2,\quad 2\leftrightarrow4,\quad 3\leftrightarrow6,\quad\dots,\quad n\leftrightarrow2n,\quad\dots
\]
\begin{itemize}
\item Different integers get different partners: if $2n=2n'$, then $n=n'$.
\item Every positive even number is used: an even number $2k$ is the partner of $k$, which is a positive integer.
\end{itemize}
So $n\mapsto2n$ is a one-to-one correspondence (a bijection) between the two sets. $\blacksquare$

For finite sets, a set always has more elements than any part of it that leaves something out. The bijection shows that this fails for infinite sets: two sets are said to have the same size (cardinality) exactly when there is a bijection between them, and by this definition the even numbers and the positive integers have the same size.
